% ---------------------------------------------------------------------------
% Pillar 2: Zero-Variance Fraction diagnostic -- eval-protocol hardening
% Berkeley harvest row 11 (validated): Yehudai et al. survey (F25 L5) +
% tau^2-Bench ablation discipline (F25 L10) applied to the iter130 per-seed
% risk-index data, plus the P2-facing slice of the Miller error-bars audit
% (row 07).
% Evidence: experiments/results/berkeley/eval_protocol_*.tsv,
%           experiments/results/berkeley/adding_error_bars_audit.tsv
% ---------------------------------------------------------------------------
\subsection{How Many Seeds Does the Headline Need? Eval-Protocol Hardening}
\label{sec:zvf-eval-protocol}

Applying the cost-efficiency lens of \citet{yehudai2025survey} and the
fine-grained-ablation discipline of $\tau^2$-Bench~\citep{barres2025tau2}
to the iter130 per-seed risk data (9 variance-mitigation methods
$\times$ 5 seeds), we ask how much of the evaluation budget the headline
claims actually require. The top-1 headline (``which method minimises
risk'') is fully seed-robust: the top-1 match rate is $1.00$ across all
$\binom{5}{k}$ seed subsets for every $k$, so the minimum viable seed
panel at the 95\% stability level is $k{=}1$ --- a five-fold evaluation
saving for that specific claim. Finer-grained claims are more expensive:
the top-3 \emph{set} reaches 80\% stability only at $k{=}4$ (top-3 match
rate $0.4$ at $k{=}1$ rising to $0.8$ at $k{=}4$), and the sign of each
method's $z$-score against \texttt{grpo} is leave-one-seed-out stable
for 8/8 non-baseline methods. Decomposing the risk index by channel,
8/9 methods are magnitude-dominant while \texttt{grpo} itself is
drift-dominant --- variance mitigation acts chiefly by suppressing the
magnitude channel --- and the best-3/mid-3/worst-3 partition is
$1.0$-stable under leave-one-seed-out re-ranking
(\texttt{experiments/results/berkeley/eval\_protocol\_summary.json}).

The earlier fused-risk AUROC and threshold analysis is withdrawn. It used
synthetic/imputed anchors and cannot validate calibrated monitoring, a safe
threshold, compute savings, or a stopping policy. The reviewed early-triage
set has only two collapsed cases, both synthetic tool use and both with early
reward zero; it provides no incremental diagnostic evidence over early reward.
Future evaluation must prospectively declare the positive class, estimand,
threshold policy, and held-out decision outcome.
